\documentclass[11pt]{article}
\usepackage[french]{babel}
\usepackage[utf8]{inputenc}
\usepackage[T1]{fontenc}
\usepackage{amsmath,amssymb,amsthm}
\usepackage{geometry}
\usepackage{hyperref}
\usepackage{listings}
\usepackage{xcolor}
\usepackage{booktabs}
\geometry{margin=2.7cm}

\newcommand{\code}[1]{\texttt{#1}}

\lstset{
  basicstyle=\ttfamily\small,
  breaklines=true,
  frame=single,
  backgroundcolor=\color{gray!8},
  numbers=none,
}

\title{LTS-RK4 multi-niveaux pour le maillage gradué en L\\
\large Journal d'implémentation et de validation}
\author{Romain Mottier}
\date{\today}

\begin{document}
\maketitle

\begin{abstract}
Ce rapport documente la conception, l'implémentation et la validation d'une
généralisation multi-niveaux du schéma explicite \emph{Local Time-Stepping}
RK4 (LTS-RK4, ``Algorithme 3'') déjà présent dans DiSk++
(\code{erk\_coupling\_hho\_scheme}), motivée par le coût prohibitif du
schéma à 2 niveaux existant sur un maillage gradué en L avec singularité de
coin rentrant. Le rapport couvre : le diagnostic du problème de coût, deux
bugs de correction trouvés dans le chemin de code ``restreint'' préexistant,
les nouvelles briques de calcul restreintes et validées, la généralisation
récursive multi-niveaux, et l'état actuel de validation.
\end{abstract}

\tableofcontents

\section{Contexte et motivation}

\subsection{Le problème de coût}

Le maillage en L gradué vers le coin rentrant (angle intérieur $\omega =
3\pi/2$) est construit par un raffinement quadtree couplé : le fond est
raffiné uniformément $\ell$ fois, et le coin est raffiné avec une
profondeur supplémentaire $\mathrm{depth} = 1.5(k+1)\,\ell$ (pour un degré
polynomial HHO $k$), afin de retrouver le taux de convergence optimal
$h^{k+1}$ malgré la singularité (théorie classique du maillage gradué).

Cette construction fait exploser le ratio $p = h_{\max}/h_{\min}$ :
$p = 1, 32, 1024, 32768,\dots$ pour les niveaux successifs $N=0,1,2,3,\dots$
(pour $k=3$, $p$ croît d'un facteur $64$ par niveau). Le schéma LTS-RK4 à 2
niveaux existant effectue $p$ sous-pas RK4 fins par macro-pas, et --- point
crucial --- \textbf{chaque appel touche l'intégralité du maillage}, pas
seulement la région ``fine''. Le coût total est donc $O(p \times
n_{\text{cells}})$, ce qui rend $N=3$ ($p\approx 32768$) de l'ordre de 9 à
10 heures de calcul, et $N=4$ totalement hors de portée.

\subsection{L'observation de Grote}

Un schéma LTS-RK4 correctement implémenté devrait coûter approximativement
la même chose qu'un RK4 classique (sans découpage coarse/fine) dès lors que
la région fine est minoritaire en nombre de degrés de liberté --- c'est
précisément ce qui n'était pas observé. Ceci a motivé une investigation en
deux temps :
\begin{enumerate}
  \item rendre le schéma à 2 niveaux \emph{effectivement} restreint en
  coût (chaque appel ne doit toucher que les DOFs de la bande active, pas
  tout le maillage) ;
  \item généraliser à un véritable schéma multi-niveaux, nécessaire car
  aucun seuil unique ne peut séparer ``fine minoritaire'' de ``coarse
  stable au grand pas de temps'' sur un maillage dont la distribution de
  tailles de cellules est quasi uniforme en échelle logarithmique sur 10 à
  15 octaves (vérifié empiriquement, voir \S\ref{sec:histogram}).
\end{enumerate}

\subsection{Diagnostic : distribution des tailles de cellules}
\label{sec:histogram}

Un histogramme des diamètres de cellules par octave ($h_{\max}/2^{k}$)
révèle que le fond uniformément raffiné (bande 0) est la seule bande
franchement majoritaire ; chaque octave suivante, jusqu'à 15 octaves de
profondeur à $N=3$, ne contient qu'un nombre quasi constant de cellules
($\sim 36$), conséquence directe de la marge de raffinement (``halo'')
ajoutée autour du coin. Cela prouve qu'\textbf{aucun seuil unique} ne peut
donner la propriété ``fine minoritaire'' de Grote sans que la bande
opposée (``coarse'') ne s'étende elle-même sur de nombreuses octaves --- ce
qui, comme démontré en \S\ref{sec:instability}, casse la validité de
l'approximation polynomiale du rôle coarse.

\section{Diagnostic empirique de deux bugs préexistants}
\label{sec:bugs}

Avant de construire quoi que ce soit de nouveau, le chemin de code
``restreint'' déjà présent dans \code{erk\_coupling\_hho\_scheme.hpp}
(\code{build\_LTS\_subspaces} + \code{erk\_weight\_LTS\_coarse\_restricted}
/ \code{erk\_weight\_LTS\_fine\_restricted}, utilisé par
\code{ERK4\_LTS\_optimised.hpp}) a été audité par lecture de code, puis
\textbf{confirmé empiriquement défaillant} par un test contrôlé : sur un
maillage localement raffiné (ratio $p=2$), avec un pas de temps
correctement résolu (CFL vérifié via une comparaison avec $p=1$, erreur
$\sim 0.02$--$0.10$, plausible), le chemin restreint donne une erreur 10 à
30 fois \emph{pire} malgré un maillage plus fin --- signature claire d'un
bug de correction, pas d'un problème de précision.

\begin{description}
  \item[Bug A --- dofs coarse jamais avancés.] Dans
  \code{erk\_weight\_LTS\_fine\_restricted}, la mise à jour RK4
  \[
    \delta x = d\tau \cdot \tfrac{1}{6}(k_1+2k_2+2k_3+k_4)
  \]
  n'est diffusée (\emph{scatter}) que sur les positions \emph{fines}
  (\code{m\_fine\_active\_c/f}) ; un commentaire affirme que les dofs
  coarses sont ``gérés ailleurs'', mais
  \code{erk\_weight\_LTS\_coarse\_restricted} ne fait \emph{jamais} muter
  l'état --- elle ne fait qu'écrire dans \code{w[]}. Aucune autre voie
  de code ne met à jour les dofs coarses : ils restent gelés à leur valeur
  initiale pendant toute la simulation.
  \item[Bug B --- terme de couplage à l'interface manquant.] Dans
  \code{compute\_w\_restricted}, la sortie ``face'' n'est calculée et
  diffusée que pour les faces \emph{classées coarses}. Or une face
  d'interface (un voisin coarse, un voisin fine) est \emph{toujours}
  classée fine par la règle d'union de \code{assemble\_P} (``une face est
  fine si au moins une cellule adjacente est fine''). La contribution de la
  cellule coarse vers cette face d'interface --- pourtant non nulle et
  physiquement nécessaire au couplage --- est donc silencieusement
  perdue.
\end{description}

Ces deux bugs affectent uniquement \code{ERK4\_LTS\_optimised.hpp} et
\code{ERK4\_LTS\_SSTAB.hpp} ; ils sont \textbf{indépendants} des
prototypes L-shape de cette session, qui utilisent l'API non restreinte
(\code{erk\_weight\_LTS\_coarse}/\code{erk\_weight\_LTS\_fine}) et restent
donc fiables. Par principe de moindre risque, ces fonctions préexistantes
n'ont \emph{pas} été modifiées ; de nouvelles fonctions, correctement
conçues dès le départ, ont été ajoutées à la place (voir \S\ref{sec:v2}).

\section{Nouvelles briques restreintes (fichier \code{erk\_coupling\_hho\_scheme.hpp})}
\label{sec:v2}

\subsection{\code{LTS\_subblock\_set} et \code{build\_LTS\_subblock}}

Contrairement à \code{build\_LTS\_subspaces} (un seul jeu de sous-blocs
coarse/fine, dans des membres de classe uniques, écrasés à chaque appel),
\code{build\_LTS\_subblock(own\_c, own\_f, halo\_rings)} retourne une
structure \emph{autonome} \code{LTS\_subblock\_set}, permettant à
plusieurs bandes de coexister --- prérequis indispensable pour le
multi-niveaux.

Le halo est construit par parcours en largeur (BFS) sur le graphe
d'adjacence cellule--face induit par la structure creuse de $K_{fc}$ :
partant de \code{own\_c}/\code{own\_f}, chaque anneau ajoute les faces
touchant une cellule déjà incluse, puis les cellules touchant une face déjà
incluse. Ce halo garantit que la chaîne d'applications répétées de
l'opérateur $B$ (utilisée par le rôle ``coarse'', voir \S\ref{sec:coarse})
reste exacte, en exploitant le fait que $K_{cc}$ est exactement
diagonale par bloc par cellule en HHO (les cellules ne couplent qu'à
travers les faces).

Chaque \code{LTS\_subblock\_set} stocke \emph{deux} jeux de sous-matrices
restreintes : \code{Kcc/Kcf/Kfc/Mc\_inv/Sff\_inv} (avec halo, pour le rôle
coarse) et leurs équivalents \code{\_own} (sans halo, pour le rôle fine, qui
n'en a pas besoin --- voir \S\ref{sec:fine}).

\subsection{Rôle ``coarse'' : \code{erk\_weight\_LTS\_coarse\_v2}}
\label{sec:coarse}

Reproduit fidèlement, mais restreint à \code{blocks.active\_c/f}
(propre + halo), l'algorithme de \code{erk\_weight\_LTS\_coarse} :
interpolation de Lagrange quadratique de la force sur $[0,dt]$, chaînes
$B^i y_n$ et $B^i F_j$ (non masquées, utilisant le halo pour rester
exactes), puis masquage \emph{final} à \code{own\_c/own\_f} uniquement
avant la dernière application de $B$ --- ce masquage final est ce qui
donne au polynôme de Taylor de sortie $w[0..3]$ sa fuite correcte vers les
faces d'interface (fixant le Bug B), sans polluer les cellules halo elles-
mêmes ($K_{cc}$ diagonale par cellule + masquage $\Rightarrow$ contribution
nulle aux lignes non-propres).

Sortie : $w_i(\tau) = w_0 + \tau w_1 + \tfrac{\tau^2}{2} w_2 +
\tfrac{\tau^3}{6} w_3$, valide en temps local sur $[0, \mathrm{len}]$ où
$\mathrm{len}$ est la durée \emph{propre} à ce niveau (voir
\S\ref{sec:mlts}, pas le macro-pas global).

\subsection{Rôle ``fine'' : \code{erk\_weight\_LTS\_fine\_v2}}
\label{sec:fine}

Un pas RK4 complet de taille $d\tau$, avec chaque étage
$k = B(P_{\text{own}} y) + w(\tau)$ --- l'entrée $y$ est masquée à
\code{own} avant application de $B$ (miroir exact de $P_{\text{fine}} \cdot
y_{\text{stage}}$ dans l'algorithme original), \emph{mais} $B$ est
appliqué avec les matrices \emph{avec halo}, de sorte que la sortie soit
non nulle non seulement sur les positions propres, mais aussi sur les
cellules coarse directement adjacentes (halo) --- exactement le
comportement de l'algorithme original non restreint, où une cellule
coarse au bord hérite d'une contribution réelle via le couplage $K_{cf}$
à la face d'interface. La mise à jour est diffusée sur \code{active\_c/f}
en entier (propre + halo), corrigeant le Bug A.

\subsection{Repli analytique : \code{erk\_weight\_LTS\_coarse\_advance\_uncovered}}

Les cellules coarse \emph{non} atteintes par le halo de la bande fine (les
cellules ``profondément intérieures'') ne reçoivent aucune correction du
rôle fine. Pour elles, $B(P_{\text{fine}} y) \equiv 0$ (aucun couplage vers
une face fine), donc leur mise à jour exacte sur tout le macro-pas est
l'intégrale fermée du polynôme cubique :
\[
  \Delta x = w_0\, dt + w_1\, \tfrac{dt^2}{2} + w_2\, \tfrac{dt^3}{6}
  + w_3\, \tfrac{dt^4}{24}.
\]
Ceci est \emph{exact}, pas une approximation : la règle de Simpson (à
laquelle se réduit la somme RK4 lorsque $B(P_{\text{fine}} y)=0$
identiquement) est exacte pour un polynôme cubique.
\textbf{Cellules couvertes par le halo et cellules non couvertes ne
doivent jamais recevoir les deux mécanismes} --- ce fut la source d'un bug
de double comptage (et d'un bug de comptage manquant dans une première
tentative), corrigé et validé (\S\ref{sec:validation}).

\section{Validation à $L=2$}
\label{sec:validation}

Méthodologie : comparaison directe, pas à pas, du vecteur d'état complet
entre l'algorithme original (non restreint, éprouvé) et les nouvelles
briques v2, sur un maillage réel avec un partage coarse/fine
\emph{authentique} (pas un cas dégénéré où l'une des deux bandes est
vide). Deux bugs ont été trouvés et corrigés durant cette validation
avant d'obtenir une correspondance exacte :

\begin{enumerate}
  \item Une première tentative (\code{fine\_v2} n'écrivant que ses
  positions propres, complétée par une intégrale fermée
  \emph{inconditionnelle} pour \emph{toutes} les cellules coarse propres)
  donnait un écart de $1.3\times 10^{-5}$ après un pas, croissant à $0.05$
  après 141 pas --- cause : les cellules coarses \emph{au bord} recevaient
  alors À LA FOIS la correction $B(P_{\text{fine}} y)$ (implicitement,
  puisque non masquée dans cette version) ET l'intégrale fermée, un double
  comptage partiel.
  \item Une seconde tentative (halo inclus dans \code{fine\_v2}, diffusion
  sur \code{active\_c/f} complet, mais \emph{sans} intégrale fermée du
  tout) donnait un écart \emph{pire} ($10^{-3}$ après un pas) --- cause :
  les cellules coarses non atteintes par le halo (environ $72/864$ dofs
  dans le cas testé) ne recevaient alors \emph{aucune} mise à jour du
  tout.
  \item La combinaison correcte --- intégrale fermée appliquée
  \emph{seulement} aux cellules propres de la bande coarse non couvertes
  par le halo de la bande fine (vérifié par appartenance d'ensemble) ---
  donne une correspondance exacte à la précision machine
  ($9.7\times 10^{-16}$, arrondi) sur les 141 pas de temps complets.
\end{enumerate}

\section{Généralisation multi-niveaux}
\label{sec:mlts}

\subsection{Structure récursive}

À chaque niveau $i < L-1$ : calculer $w_i$ (rôle coarse, avec la durée
\emph{propre} à ce niveau $\mathrm{len}_i$, pas le macro-pas global),
combiner avec la contribution ancêtre reçue en argument (somme
coefficient par coefficient, les deux étant exprimés dans le même repère
temporel local), avancer les cellules propres non couvertes via
l'intégrale fermée du polynôme \emph{combiné}, puis boucler $p_i$ fois :
à chaque sous-intervalle, décaler (\emph{Taylor-shift}) le polynôme
combiné vers la nouvelle origine locale, et soit appeler
\code{erk\_weight\_LTS\_fine\_v2} directement (cas de base, $i+1=L-1$),
soit récurser.

\paragraph{Décalage de Taylor.} Pour un polynôme cubique $T(\tau) = w_0 +
w_1\tau + \tfrac{w_2}{2}\tau^2 + \tfrac{w_3}{6}\tau^3$, la ré-expansion
exacte autour d'un nouveau point $a$ est :
\[
  w_0' = T(a),\quad w_1' = T'(a),\quad w_2' = T''(a),\quad w_3' = w_3.
\]
Ceci est une identité exacte (série de Taylor finie d'un polynôme de degré
3), pas une approximation.

\subsection{Point clé : mise à l'échelle du pas de temps par niveau}
\label{sec:instability}

Une première tentative de \S\ref{sec:mlts} appliquée à un simple split à
2 niveaux, en réduisant agressivement le seuil $h_c$ (pour minimiser la
région ``fine'') \emph{tout en gardant le même pas de temps macro
$dt=\mathrm{cfl}\cdot h_{\max}$}, a produit un résultat \code{NaN} ---
même avec $h_c = 16\,h_{\min}$ (bande coarse s'étalant encore sur $\sim 6$
octaves). Diagnostic : le polynôme cubique $w[]$ n'a que 4 degrés de
liberté ; sur une fenêtre $dt$ calée pour des cellules de taille
$h_{\max}$, il ne peut pas représenter la dynamique de cellules bien plus
petites oscillant plusieurs fois pendant cette même fenêtre --- rupture de
validité inhérente à la méthode, pas un bug d'implémentation (les
coefficients eux-mêmes restent corrects, comme vérifié en
\S\ref{sec:validation}).

\textbf{Correction} : chaque niveau $i$ reçoit son propre $dt_i =
\mathrm{cfl} \cdot \mathrm{scale}_i$, où $\mathrm{scale}_i$ est la taille
caractéristique des cellules de ce niveau ($h_{\max}$ pour le niveau 0,
le seuil $h_{c,i}$ pour les niveaux intermédiaires, $h_{\min}$ pour le
niveau le plus fin). Le ratio de branchement est $p_i = \mathrm{round}
(dt_i / dt_{i+1})$.

\subsection{Validation de la récursion générique}

Le pilote récursif générique, spécialisé à $L=2$ avec exactement le même
seuil que la validation de \S\ref{sec:validation}, reproduit le résultat
de référence à la précision machine ($6.9\times 10^{-16}$). Ceci confirme
que la récursion elle-même (décalage de Taylor, accumulation des
contributions ancêtres, repli analytique) est correcte --- tout écart
constaté à $L>2$ est donc une question de réglage (précision), pas de
correction.

\subsection{Résultats à $L=5$ ($N=2$, $k=3$)}

Avec des bandes espacées d'environ 2 octaves chacune (5 niveaux au total
pour couvrir les 10 octaves de $p=1024$) :
\begin{itemize}
  \item Aucune divergence (\code{NaN}) --- confirmation de la correction
  de mise à l'échelle du pas de temps.
  \item Temps de calcul : $170\,\mathrm{s}$, contre $\sim 540\,\mathrm{s}$
  pour le schéma original non restreint --- gain $\times 3{,}2$.
  \item Précision : erreur $L^2 = 3{,}7\times 10^{-3}$, contre
  $1{,}14\times 10^{-5}$ pour le schéma de référence --- une dégradation
  significative (facteur $\sim 326$), attribuable à l'accumulation
  d'approximations cubiques à travers 4 niveaux intermédiaires plutôt qu'à
  une erreur de conception.
\end{itemize}

\section{Bugs supplémentaires trouvés à $L=3$ et au-delà}
\label{sec:l3bugs}

Suite à la validation de la récursion générique à $L=2$ (\S\ref{sec:mlts}),
un test contrôlé similaire à $L=3$ (comparaison contre le schéma à 2
niveaux original avec $P_{\text{fine}}=\{\text{bande}_1+\text{bande}_2\}$,
sur un maillage réduit $N=1$, $p_{\text{global}}=32$, pour qu'un vrai bug
se manifeste de façon flagrante plutôt que d'être masqué par une erreur
d'approximation légitime) a révélé un écart de \textbf{facteur 109} par
rapport à la référence --- bien trop grand pour être expliqué par
l'approximation supplémentaire d'un unique niveau intermédiaire.

\paragraph{Bug C --- mauvaise cible de couverture.} La fonction
\code{erk\_weight\_LTS\_coarse\_advance\_uncovered} vérifiait, pour la
bande $i$, la couverture par le halo de la bande $i{+}1$
(\code{blocks[i+1]}). Ceci est correct \emph{uniquement} quand $i{+}1$ est
la bande terminale (celle qui reçoit réellement \code{erk\_weight\_LTS\_-
fine\_v2}, donc une écriture directe). Pour toute bande intermédiaire
($i+1 < L-1$), son propre halo ne sert que de contexte en lecture seule
pour son \emph{propre} calcul de $B$-chaîne --- elle ne diffuse jamais
d'écriture vers les cellules qui l'y ont atteinte. Le correctif : la
vérification de couverture doit \emph{toujours} référencer la bande
terminale \code{blocks[L-1]}, pour tous les niveaux $i$, pas la bande
$i{+}1$ immédiate. Ce correctif a fait chuter le ratio d'erreur de
$\times 109$ à $\times 17$.

\paragraph{Bug D --- désaccord de quantification $p_0 \times p_1 \times
\dots \ne p_{\text{global}}$.} Les ratios $p_i$ de chaque niveau étaient
arrondis indépendamment (\code{round(dt\_i/dt\_{i+1})}), sans garantir que
leur produit égale exactement $p_{\text{global}}$ --- introduisant un
écart artificiel de résolution temporelle fine entre le schéma récursif et
la référence. En choisissant les seuils pour que le produit soit exact
(ex. $p_0=4, p_1=8$ pour $p_{\text{global}}=32$), le ratio est retombé à
$\times 8{,}7$ --- jugé plausible (pas de bug flagrant restant à ce point
de test précis).

\subsection{Le ratio d'erreur croît avec $L$, indépendamment de $p_{\text{global}}$}

Après application du correctif du Bug C au pilote générique complet
(\code{ERK4\_MLTS\_timing\_test.hpp}), une série de tests contrôlés
montre :
\begin{center}
\begin{tabular}{lccc}
\toprule
Configuration & $L$ & $p_{\text{global}}$ & ratio erreur/référence \\
\midrule
$N=1$, bandes larges (Bug D corrigé) & 3 & 32 & $\times 8{,}7$ \\
$N=1$, bandes étroites ($\sim$1 octave) & 5 & 32 & $\times 20{,}1$ \\
$N=2$, bandes larges ($\sim$2 octaves) & 5 & 1024 & $\times 92$ \\
\bottomrule
\end{tabular}
\end{center}

Le ratio croît \emph{à la fois} avec $L$ (à $p_{\text{global}}$ fixé, 8,7
$\to$ 20,1 en ajoutant 2 niveaux) et avec $p_{\text{global}}$ (à $L$ fixé,
20,1 $\to$ 92 en passant de 32 à 1024). Ce n'est plus la signature d'un
bug isolé et localisable comme les Bugs A--D : c'est un pattern cohérent
de dégradation qui ressemble à une limite \emph{inhérente} de
l'empilement de plusieurs approximations cubiques indépendantes (chaque
niveau intermédiaire introduit sa propre erreur de troncature, et ces
erreurs semblent se composer plutôt que simplement s'additionner). Cette
hypothèse n'est pas encore prouvée rigoureusement --- elle reste la
piste la plus probable au vu des données recueillies, mais d'autres bugs
plus subtils ne peuvent pas être formellement exclus sans une
investigation plus poussée.

\section{Correction structurelle majeure : chaque niveau reçoit une vraie dynamique}
\label{sec:realrk4}

L'hypothèse de \S\ref{sec:l3bugs} (limite inhérente de l'empilement
d'approximations cubiques) s'est révélée \textbf{incorrecte} : le
problème était une erreur de conception structurelle, identifiée grâce à
une clarification de l'utilisateur sur le fonctionnement voulu du
multi-niveaux.

\paragraph{Le design fautif.} Dans la version précédente, \emph{seul le
niveau terminal} (le plus fin) recevait un véritable pas RK4
(\code{erk\_weight\_LTS\_fine\_v2}) ; tous les niveaux intermédiaires
n'avançaient leurs propres cellules que via l'intégrale fermée
(analytique) du polynôme de Taylor. Ceci sous-estime systématiquement la
dynamique réelle de tous les niveaux sauf les deux extrêmes.

\paragraph{Le design corrigé.} Le niveau 0 (le plus grossier) reste
purement analytique (jamais de RK4 réel pour lui-même) --- mais
\textbf{tous les niveaux $i\geq1$}, pas seulement le dernier, reçoivent
un véritable pas RK4 (\code{erk\_weight\_LTS\_fine\_v2}), \emph{un seul
pas} de taille $\mathrm{len}$ (l'intervalle propre à ce niveau, qui par
construction égale $dt_{\text{level}(i)}$), en utilisant le polynôme
figé $w_{i-1}$ du niveau parent comme terme de forçage additif. Chaque
niveau $i<L-1$ calcule en plus son propre $w_i$ (\emph{avant} que son
propre pas RK4 ne modifie l'état, pour respecter la convention ``figé au
début de l'intervalle'') afin de le transmettre au niveau $i+1$. Le halo
de protection (utilisé à la fois pour la précision de la chaîne $B$ du
rôle coarse et pour la diffusion réelle du rôle fine vers la couche
adjacente $p{-}1$) a aussi été réduit de 6 à 3 anneaux, plus proche de
l'idée d'une simple ``couche de protection'' plutôt qu'une portée
générique large.

\paragraph{Résultat.} Validé d'abord sur le cas contrôlé $L=3$, $N=1$,
$p_{\text{global}}=32$ : le ratio d'erreur par rapport à la référence
chute de $\times 8{,}7$ à $\times 1{,}22$ --- un résultat quasi exact.
Propagé au pilote général et testé à l'échelle problématique ($L=5$,
$N=2$, $p_{\text{global}}=1024$, celle qui donnait $\times 92$
auparavant) :
\begin{itemize}
  \item Ratio d'erreur : $\times 9{,}1$ (contre $\times 92$) --- une
  réduction d'un facteur $\sim 10$.
  \item Temps de calcul : $142\,\mathrm{s}$ --- \emph{plus rapide} que la
  version précédente ($165\,\mathrm{s}$), malgré davantage de travail par
  niveau, grâce au halo réduit.
  \item Comparé au schéma original non restreint ($\sim 540\,\mathrm{s}$) :
  gain de vitesse $\times 3{,}8$, pour une précision maintenant dans une
  plage raisonnable pour un algorithme multi-niveaux authentiquement
  différent (pas une preuve de bug supplémentaire).
\end{itemize}

\section{L'erreur résiduelle croît avec $L$, pas avec $p_{\text{global}}$}
\label{sec:lgrowth}

Après la correction structurelle de la \S\ref{sec:realrk4}, un résidu
d'erreur subsistait : $\times 1{,}22$ à $L=3$ mais $\times 9{,}1$ à
$L=5$. L'utilisateur a formellement rejeté l'idée qu'un design
multi-niveaux correctement implémenté puisse introduire une erreur non
nulle, et a demandé d'isoler la cause précise plutôt que d'accepter ce
résidu comme une limite inhérente.

\paragraph{Isolation de la variable.} Un test contrôlé à $L=5$,
$N=1$, $p_{\text{global}}=32$ (petit, donc rapide, mais avec le même
nombre de niveaux que le cas problématique) a donné un ratio de
$\times 9{,}08$ --- quasi identique au $\times 9{,}1$ obtenu à
$p_{\text{global}}=1024$. Ceci démontre que la croissance de l'erreur
suit \textbf{le nombre de niveaux $L$, pas le rapport de raffinement
global}.

\paragraph{Deux hypothèses testées et écartées.}
\begin{enumerate}
  \item \emph{Terme de forçage ancestral manquant.} L'hypothèse que
  \code{erk\_weight\_LTS\_coarse\_v2} d'un niveau intermédiaire ignorait
  la contribution de son propre ancêtre (au lieu de la geler comme un
  instantané) a été corrigée proprement (nouveau paramètre
  \code{ancestor\_w}, incorporé sans réappliquer $M_c^{-1}$ deux fois ---
  une première tentative naïve avait fait exploser le calcul par erreur
  d'unités). Effet mesuré~: $\times 9{,}08 \to \times 9{,}02$ ---
  négligeable.
  \item \emph{Halo trop étroit.} Élargir \code{halo\_rings} de 3 à 10
  \textbf{a empiré} le résultat ($\times 21$ au lieu de $\times 3$ à
  $L=3$), écartant complètement cette piste.
\end{enumerate}

\paragraph{Un balayage précis en $L$, à $p_{\text{global}}=32$ fixe,
révèle un saut net et un plateau.} En gardant \code{band0} identique
(même seuil $h_{\max}/2$) pour tout $L$ et en variant seulement le
nombre de bandes intermédiaires :
\begin{center}
\begin{tabular}{@{}lccccc@{}}
\toprule
$L$ & 2 & 3 & 4 & 5 \\
\midrule
Ratio & $1{,}00$ & $2{,}97$ & $9{,}00$ & $9{,}02$ \\
\bottomrule
\end{tabular}
\end{center}
Le saut se produit précisément entre $L=3$ et $L=4$, puis
\textbf{plafonne}. La différence structurelle entre $L=3$ et $L=4$ :
c'est la première fois qu'un niveau ``rôle coarse'' reçoit son
forçage ancestral d'un \emph{autre} niveau lui-même en rôle coarse
(chaîne \code{band1}$\to$\code{band2}$\to$\code{band3}), plutôt que de
recevoir directement de la racine ou de nourrir directement le niveau
terminal. Un test à un seul macro-pas (\code{nt=1}) montre par ailleurs
que l'erreur \emph{par pas} reste minuscule ($\sim 10^{-4}$ relatif) à
tous les $L$ --- ce n'est donc pas un bug isolé dans un nœud de la
récursion, mais une erreur qui \emph{s'amplifie} sur les 141 pas,
spécifiquement quand une telle chaîne existe.

\section{Relecture de l'article de Diaz--Grote : la cause profonde}
\label{sec:groteread}

Sur suggestion de l'utilisateur, l'article fondateur
\emph{Multi-Level Explicit Local Time-Stepping Methods for Second-Order
Wave Equations} (Diaz \& Grote, CMAME 2015) et sa version 2-niveaux
(Grote \& Mitkova, 2012) ont été lus en détail (PDF fournis par
l'utilisateur après plusieurs tentatives infructueuses d'accès en
ligne). L'Algorithme~4 de l'article (MLTS2, récursion multi-niveaux)
révèle une différence structurelle fondamentale avec notre
implémentation~:
\begin{quote}
\code{MLTS2(y\_inter, w, $\Delta t$, l)}: si $l<N_{level}$, appelle
récursivement \code{MLTS2(y\_inter, w -- A($P_l$--$P_{l+1}$)y\_inter,
$\Delta t/p_l$, l+1)}, où le terme \code{A($P_l$--$P_{l+1}$)y\_inter}
est \textbf{recalculé à chaque sous-pas}, à partir de l'état
\code{y\_inter} \emph{réel et à jour} --- jamais extrapolé.
\end{quote}
Deux différences architecturales expliquent l'écart observé :
\begin{enumerate}
  \item \textbf{Extrapolation vs recalcul exact.} Notre design calculait
  \code{own\_w\_i} \emph{une fois} par visite d'un niveau, puis
  l'extrapolait analytiquement (\emph{Taylor-shift}) sur tous les
  sous-pas du niveau suivant. Grote--Diaz recalculent cette contribution
  \emph{à chaque sous-pas}, à partir de l'état réel --- pas une
  approximation polynomiale figée sur tout l'intervalle macro.
  \item \textbf{Opérateur global vs sous-blocs restreints.} Tout
  l'algorithme de Grote--Diaz opère sur le vecteur et la matrice
  \emph{globaux} (masqués via des projecteurs diagonaux $P_l$), jamais
  sur des sous-blocs restreints à un voisinage local. La portée de la
  ``fuite'' d'un niveau vers ses voisins (via les puissances
  $(AP)^j$ dans la preuve, Lemme~4.1) s'étend avec le nombre de
  sous-pas $p$, pas avec une largeur de halo fixe --- ce qui explique
  \emph{a posteriori} pourquoi élargir notre halo restreint (fixe) a
  échoué.
\end{enumerate}
Une première tentative d'adopter fidèlement cette structure (niveau
terminal seul avec vraie dynamique RK4, niveaux intermédiaires en
``pass-through'') a révélé un piège supplémentaire~: sans le mécanisme
de fuite propre à chaque niveau intermédiaire, les cellules de la
racine adjacentes à la première bande fine restaient orphelines (jamais
mises à jour), faisant \emph{exploser} l'erreur (band0~: $0{,}48$,
contre $1{,}7\times10^{-5}$ dans le design original). Plusieurs correctifs
ciblés (recalcul de la racine à chaque sous-pas, cible de couverture
élargie) ont réduit mais pas éliminé ce résidu ($\sim\times20$,
plafonnant avec $L$ mais toujours pire qu'à l'origine) --- confirmant
que rester sur des sous-blocs restreints ne permet pas de reproduire
fidèlement le mécanisme de l'article.

\section{Résolution : réécriture fidèle sur matrices globales}
\label{sec:globalfix}

Sur décision explicite de l'utilisateur (``il faut prendre leur
design''), le pilote multi-niveaux a été entièrement réécrit pour
suivre l'Algorithme~4 de Grote--Diaz \emph{littéralement}, tout en
gardant RK4 (et non le leap-frog de l'article) comme intégrateur de
base, conformément à l'algorithme déjà validé (Algorithme~3) de DiSk++.

\paragraph{Nouvelle brique : \code{erk\_weight\_LTS\_coarse\_global}.}
La fonction globale existante \code{erk\_weight\_LTS\_coarse} met en
cache sa liste de degrés de liberté actifs au premier appel
(\code{m\_coarse\_active\_dofs}) et la réutilise indéfiniment --- sûr
pour ses 11 appelants existants (qui n'utilisent jamais qu'un seul
\code{Pcoarse} fixe), mais incompatible avec un pilote multi-niveaux
appelant la fonction avec un \code{Pcoarse} différent à chaque bande et
à chaque sous-pas. Une nouvelle fonction, strictement additive et sans
état (\code{erk\_weight\_LTS\_coarse\_global}), reproduit exactement le
même calcul mais reconstruit sa liste de degrés de liberté actifs à
chaque appel.

\paragraph{Nouveau pilote récursif.} Pour chaque niveau $i<L-1$, à
chaque sous-pas court $m$~:
\begin{enumerate}
  \item recalculer \emph{à neuf}, depuis l'état courant \code{xn}, la
  contribution propre de la bande $i$ via
  \code{erk\_weight\_LTS\_coarse\_global(xn, $P_{\text{band}_i}$,
  ...)} --- avec $P_{\text{band}_i}$ un projecteur \emph{global}
  (diagonal, taille du maillage complet), pas un sous-bloc restreint ;
  \item \textbf{additionner} cette contribution à celle héritée des
  ancêtres (\emph{Taylor-shiftée} à l'origine locale de ce sous-pas)
  --- exactement le \code{v := w + $\sum_k$ w\_k} de l'article ;
  \item récurser vers le niveau $i+1$ avec cette somme.
\end{enumerate}
Au niveau terminal $L-1$, un unique pas RK4 authentique
(\code{erk\_weight\_LTS\_fine}, la fonction globale déjà validée,
inchangée) est appliqué, forcé par la somme complète accumulée depuis
la racine. \textbf{Aucune} étape séparée de type ``avancer les degrés
de liberté non couverts'' n'est nécessaire : chaque niveau, y compris
la racine, est mis à jour uniquement par la fuite naturelle des
opérateurs globaux, enchaînée à travers les nombreuses visites
imbriquées de la récursion --- exactement comme dans l'article, où seul
l'appel de base met directement à jour le vecteur d'état.

\paragraph{Résultat : validation exacte à tous les niveaux testés.}
Sur le même test contrôlé ($N=1$, $p_{\text{global}}=32$)~:
\begin{center}
\begin{tabular}{@{}lcccc@{}}
\toprule
$L$ & 2 & 3 & 4 & 5 \\
\midrule
Ratio & $1{,}0000000000$ & $1{,}0000000000$ & $0{,}9999999999$ & $0{,}9999999999$ \\
Erreur max & $\sim 4\times10^{-11}$ & $\sim 4\times10^{-11}$ & $\sim 4\times10^{-11}$ & $\sim 6\times10^{-11}$ \\
\bottomrule
\end{tabular}
\end{center}
Précision machine, \textbf{sans aucune dégradation avec $L$} --- pour la
première fois depuis le début de cette investigation. Ceci confirme que
l'écart observé n'était pas une limite inhérente au multi-niveaux, mais
bien un choix architectural incorrect (extrapolation figée sur des
sous-blocs restreints) qui ne reproduisait pas fidèlement le mécanisme
prouvé stable et exact de l'article.

\section{État actuel et travail restant}

\begin{itemize}
  \item \textbf{Correction fondamentale validée} : le pilote
  multi-niveaux réécrit fidèlement sur matrices globales
  (\S\ref{sec:globalfix}) reproduit la référence à précision machine
  pour $L=2,3,4,5$, sans dégradation avec la profondeur --- résolvant
  définitivement le problème identifié en \S\ref{sec:lgrowth}.
  \item \textbf{Quatre bugs réels trouvés et corrigés en tout} : mauvaise
  cible de couverture (Bug C), désaccord de quantification (Bug D),
  seul le niveau terminal recevait une vraie dynamique RK4
  (\S\ref{sec:realrk4}), et le mécanisme d'extrapolation figée par
  sous-bloc au lieu du recalcul exact sur matrices globales
  (\S\ref{sec:globalfix}).
  \item \textbf{Coût actuel non optimisé} : la version corrigée
  utilise les matrices \emph{globales} (produits matrice-vecteur creux
  sur tout le maillage) à chaque sous-pas, recalculés très fréquemment
  --- plus coûteux que la version (incorrecte) à sous-blocs restreints.
  L'algorithme est maintenant \emph{correct} ; le travail restant est de
  retrouver un coût réduit (probablement en revenant à des sous-blocs
  restreints, mais construits correctement cette fois, maintenant que
  l'algorithme de référence est connu) pour que le multi-niveaux tienne
  sa promesse de réduction de coût à l'échelle $N=3$/$N=4$.
  \item \textbf{Prochaine étape} : optimiser le coût, puis appliquer ce
  design corrigé au balayage complet $N=0$ à $N=4$ pour obtenir la
  courbe de convergence à 5 points pour $k=3$.
\end{itemize}

\section{Tentatives d'optimisation restreinte : trois pistes explorées et écartées}
\label{sec:restrictedattempts}

Suite à la demande de l'utilisateur (\og il faut réussir à optimiser
l'approche de Grote qui est exacte \fg{}), plusieurs tentatives ont été
faites pour retrouver le bénéfice de performance des sous-blocs
restreints (\S\ref{sec:realrk4}) \emph{sans} réintroduire l'erreur
dépendante de $L$, en réutilisant le mécanisme exact de
\S\ref{sec:globalfix} :

\begin{enumerate}
  \item \textbf{Couverture élargie + \code{advance\_uncovered}.} Chaque
  niveau intermédiaire calcule sa contribution propre sur un sous-bloc
  restreint, puis l'applique directement à ses propres degrés de liberté
  via une intégrale fermée (au lieu de compter sur la fuite du niveau
  terminal, dont le halo restreint ne peut jamais atteindre band0).
  Résultat~: la magnitude totale du déplacement de band0 devient
  \emph{exacte} ($0{,}14497712$, identique à 6 chiffres près à la
  version globale), mais l'erreur \emph{résiduelle} reste à
  $\sim\times20$, stable avec $L$.
  \item \textbf{Approche hybride} (RK4 restreint \emph{et} intégrale
  fermée combinés pour chaque niveau)~: rejetée sur instruction de
  l'utilisateur avant même d'être testée à fond, car elle mélangeait
  deux mécanismes sans justification théorique claire -- un test rapide
  a effectivement dégradé $L=2$ (qui était pourtant exact) à
  $\times16{,}8$.
  \item \textbf{Rafraîchissement à la fréquence du niveau terminal.}
  Hypothèse~: band0 n'est recalculé qu'une fois par son propre intervalle
  grossier, alors que dans la version globale son influence est
  implicitement réévaluée à chaque sous-pas du niveau terminal (bien plus
  fréquent). Le code a été restructuré pour recalculer et appliquer la
  contribution de \emph{chaque} niveau ancêtre à la granularité fine du
  terminal. Résultat~: \textbf{aucune amélioration mesurable}
  ($\times19{,}80$ contre $\times19{,}81$ avant) -- réfutant cette
  hypothèse spécifique.
\end{enumerate}

La cause exacte du résidu $\times20$ (magnitude correcte, distribution
légèrement fausse) reste donc \textbf{non résolue} à l'issue de cette
session. Sur décision de l'utilisateur, la suite privilégie
l'utilisation de la version \emph{globale exacte} (\S\ref{sec:globalfix})
pour produire des résultats de confiance, en acceptant son coût plus
élevé, plutôt que de continuer à chercher une optimisation restreinte
sans nouvelle piste solide.

\section{Résultats de convergence obtenus avec la version globale exacte}
\label{sec:finalresults}

Le balayage $N=0,1,2$ a été exécuté avec la version globale exacte
(\S\ref{sec:globalfix}), chaque appel du niveau terminal et de chaque
niveau ancêtre utilisant les matrices globales $K_{cc}, K_{cf}, K_{fc}$
non restreintes.

\begin{center}
\begin{tabular}{@{}lccccc@{}}
\toprule
$N$ & $h_{\max}$ & $p_{\text{global}}$ & $L$ & Erreur $L_2$ & Temps (s) \\
\midrule
0 & 0,3536 & 1     & 1 & $2{,}193\times10^{-2}$ & 0,008 \\
1 & 0,1768 & 32    & 2 & $2{,}769\times10^{-4}$ & 13,7 \\
2 & 0,0884 & 1024  & 5 & $1{,}145\times10^{-5}$ & 2331 ($\approx$39 min) \\
\bottomrule
\end{tabular}
\end{center}

\textbf{Validation croisée cruciale}~: l'erreur $L_2$ obtenue à $N=2$
($1{,}144980932\times10^{-5}$) correspond, à 9 chiffres significatifs
près, à la valeur de référence de confiance obtenue en tout début de
session avec le schéma 2-niveaux non restreint
($1{,}14498093168055\times10^{-5}$) -- confirmation indépendante et
définitive que l'algorithme multi-niveaux (avec la correction de
\S\ref{sec:globalfix}) reproduit exactement la physique attendue, y
compris à l'échelle réelle du maillage ($N=2$, $p_{\text{global}}=1024$,
pas seulement le petit cas contrôlé $N=1$).

\textbf{Ordre de convergence observé}~:
\[
  \text{ordre}(N{=}0{\to}1) = \log_2\!\frac{2{,}193\times10^{-2}}{2{,}769\times10^{-4}} \approx 6{,}31,
  \qquad
  \text{ordre}(N{=}1{\to}2) = \log_2\!\frac{2{,}769\times10^{-4}}{1{,}145\times10^{-5}} \approx 4{,}60.
\]
L'ordre $N{=}1{\to}2$ ($\approx4{,}60$) correspond exactement à la
valeur trouvée en tout début de session avec le schéma de confiance
non restreint, confirmant que $k=3$ atteint un ordre quasi-optimal
(l'ordre optimal théorique étant $k+1=4$; la légère sur-convergence
provient du maillage gradué en coin).

\textbf{Coût et arrêt à $N=2$}~: le calcul de $N=3$
($p_{\text{global}}=32768$, $32\times$ plus grand que $N=2$) a été
lancé puis interrompu après extrapolation du coût observé
($\text{nt}\times p_{\text{global}}\times n_{\text{cells}}$), estimé à
environ \textbf{4 jours} avec la version globale non optimisée --
jugé impraticable pour cette session. Sur décision de l'utilisateur,
la session s'arrête avec 3 points de courbe de convergence
($N=0,1,2$), tous confirmés corrects, plutôt que de laisser tourner un
calcul de plusieurs jours sans supervision.

\section{Bilan et travail restant}

\begin{itemize}
  \item \textbf{Acquis définitif}~: l'algorithme multi-niveaux
  fidèle à Grote--Diaz (\S\ref{sec:globalfix}) est \emph{mathématiquement
  correct}, validé à la fois par précision machine sur cas contrôlés
  ($L=2..5$) et par une correspondance à 9 chiffres avec la référence
  de confiance sur un cas réel ($N=2$).
  \item \textbf{Non résolu}~: l'optimisation vers des sous-blocs
  restreints réintroduit un résidu $\sim\times20$ dont la cause précise
  n'a pas été identifiée malgré trois tentatives motivées
  (\S\ref{sec:restrictedattempts}). Cette piste reste ouverte pour une
  prochaine session.
  \item \textbf{Résultat scientifique obtenu}~: 3 points de convergence
  ($N=0,1,2$) pour $k=3$, confirmant un ordre quasi-optimal
  ($\approx4{,}60$), cohérent avec les résultats obtenus en début de
  session par une méthode différente (schéma 2-niveaux non restreint).
  \item \textbf{Prochaine étape}~: soit reprendre l'optimisation
  restreinte (pour rendre $N=3$/$N=4$ praticables), soit accepter le
  coût de la version globale exacte pour un nombre limité de points
  supplémentaires avec un budget de temps dédié plus important.
\end{itemize}

\section{Recalibrage CFL : un gain de performance confirmé et indépendant}
\label{sec:cflfix}

Une piste totalement orthogonale à l'architecture multi-niveaux a été
explorée et validée : le facteur de sécurité CFL utilisé pour fixer le
pas de temps macro (\texttt{cfl\_factor}, dans \texttt{dt\_macro =
cfl\_factor} $\times\, h_{\max}$) était hérité d'un choix ancien et
jamais recalibré. Une recherche binaire empirique de la limite de
stabilité RK4 réelle (à $N=1$) a montré que la vraie limite se situe
vers \texttt{cfl\_factor}$\approx 0{,}08$--$0{,}085$, alors que la valeur
utilisée était $0{,}002$ -- soit un pas de temps $\sim$40$\times$ plus
petit que nécessaire, sans aucun bénéfice de précision (RK4 est déjà
d'ordre 4 ; l'erreur spatiale domine largement l'erreur temporelle dans
cette plage). Avec \texttt{cfl\_factor}$=0{,}05$ (marge de sécurité
sous la limite réelle), $N=2$ passe de 2330\,s à 96--111\,s
($\sim$25$\times$ plus rapide), avec une erreur $L_2$ identique à 6
chiffres significatifs près par rapport à \texttt{cfl\_factor}$=0{,}002$.
Ce gain est \emph{confirmé, indépendant de l'exactitude de
l'algorithme}, et s'applique tel quel au balayage $N=0..4$.

\section{Quatre tentatives d'optimisation algorithmique supplémentaires}
\label{sec:moreattempts}

Suite à la demande explicite de continuer à chercher une optimisation
\emph{algorithmique} (plutôt que de changer de technologie de
parallélisation -- voir \S\ref{sec:parallel}), quatre pistes
supplémentaires ont été explorées ce soir, en s'appuyant notamment sur
une relecture attentive d'Almquist \& Mehlin (\emph{Multi-level local
time-stepping methods of Runge-Kutta type for wave equations}, CRC
Preprint 2016/16) -- l'extension RK explicite du schéma leapfrog de
Diaz--Grote, directement pertinente puisque notre schéma utilise RK4.

\begin{enumerate}
  \item \textbf{Terme croisé ancestor\_w} (\S\ref{sec:restrictedattempts}
  poursuivi) : leur formule (22) repropage explicitement la contribution
  de chaque niveau ancêtre \emph{à travers} l'opérateur $B P_l$ du niveau
  courant, plutôt que de simplement l'additionner après décalage de
  Taylor. Un mécanisme équivalent existait déjà dans le code
  (\texttt{ancestor\_w}) mais n'était pas branché dans le driver restreint
  le plus récent. Rebranché et testé sur $N=2$ réel : résultat
  \textbf{pire} ($1{,}92\times10^{-4}$ contre $1{,}54\times10^{-4}$ sans ce
  terme, référence $1{,}14\times10^{-5}$) -- très probablement un
  double-comptage avec le mécanisme \texttt{advance\_uncovered\_full}
  déjà en place. Abandonné.
  \item \textbf{Micro-optimisations sans risque} : (a) un calcul
  gaspillé identifié dans \texttt{erk\_weight\_LTS\_coarse\_global} (la
  partie face de \texttt{B\_zero\_y(IPF\_k)} n'est jamais lue, seule la
  partie cellule sert) ; (b) le scan $O(n_{\text{dof}})$ de la diagonale
  \texttt{Pcoarse} pour reconstruire la liste des degrés de liberté
  actifs, refait à \emph{chaque appel} alors qu'elle est identique pour
  toute la simulation à un niveau donné -- remplacé par une liste
  précalculée passée directement. Les deux corrections sont
  \textbf{validées correctes} (résultats identiques bit à bit sur $N=1$
  et $N=2$) mais \textbf{n'apportent aucun gain de temps mesurable}
  ($110{,}9$\,s contre $96$--$111$\,s avant, à $N=2$) -- confirmant que
  le coût dominant est le \emph{volume brut} de produits
  matrice-vecteur creux exigé par l'algorithme, pas un gaspillage
  identifiable.
  \item \textbf{Parallélisation OpenMP des produits matrice-vecteur
  creux} (\S\ref{sec:parallel}) : tentative directe (une région
  parallèle OpenMP par appel), mesurée \textbf{beaucoup plus lente}
  ($N{=}1$ : $18$--$26$\,s contre $0{,}67$\,s en série) -- le coût de
  synchronisation par appel dépasse largement le travail à paralléliser,
  cette fonction étant appelée des dizaines de milliers de fois avec un
  travail minuscule à chaque fois. Abandonnée (code source retiré,
  gardé en commentaire).
  \item \textbf{Hybride : ancêtres exacts, terminal restreint seul} :
  hypothèse que le niveau terminal (visité $\propto p_{\text{global}}$
  fois, donc le principal moteur de coût) pourrait être restreint à un
  sous-bloc local (\texttt{erk\_weight\_LTS\_fine\_v2}) tout en gardant
  les niveaux ancêtres en calcul global exact (peu visités, donc leur
  coût $O(n_{\text{dof}})$ reste supportable), avec une avance directe en
  forme fermée des degrés de liberté propres de chaque ancêtre (utilisant
  cette fois le \texttt{own\_w\_i} \emph{exact}, pas l'approximation
  restreinte). Testé sur $N=2$ réel : \textbf{pire sur les deux plans} --
  précision dégradée ($1{,}09\times10^{-3}$, near $100\times$ la
  référence) \emph{et} aucun gain de vitesse mesurable ($88{,}8$\,s,
  comparable à la version globale complète). Ceci réfute l'hypothèse que
  le niveau terminal domine seul le coût total -- les niveaux ancêtres,
  cumulés sur toute la profondeur de récursion, pèsent visiblement tout
  autant. Abandonné.
\end{enumerate}

\textbf{Bilan honnête} : quatre tentatives de restructuration
algorithmique n'ont amélioré ni la précision ni la vitesse par rapport
à la version globale exacte déjà validée. Les deux micro-optimisations
sans risque sont correctes mais négligeables en pratique. Le
recalibrage CFL (\S\ref{sec:cflfix}) reste le seul gain de performance
substantiel et confirmé obtenu ce soir, complètement indépendant de
cette question.

\section{Parallélisation : pourquoi elle n'aide pas ici}
\label{sec:parallel}

La question a été posée d'utiliser OpenMP, MPI, ou un GPU pour
accélérer l'algorithme exact. Le diagnostic est le même dans les trois
cas et a été confirmé empiriquement pour OpenMP (point 3 ci-dessus) :
l'algorithme est une longue chaîne d'opérations \emph{séquentielles},
petites et bon marché (chaque étage RK4 dépend du précédent, chaque
niveau de récursion dépend du précédent, chaque pas de temps dépend du
précédent). Le seul parallélisme disponible sans restructuration
profonde se trouve \emph{à l'intérieur} de chaque produit
matrice-vecteur creux -- mais ces produits sont appelés des dizaines de
milliers de fois avec un travail minuscule à chaque fois, et le
surcoût de synchronisation (mémoire partagée pour OpenMP, réseau/IPC
pour MPI, lancement de kernel pour GPU -- dans cet ordre croissant de
coût) domine systématiquement le travail utile. Un vrai gain
nécessiterait de regrouper de nombreuses petites opérations en moins
d'opérations plus grosses (tâches OpenMP à gros grain, CUDA Graphs) --
un travail de conception substantiel, hors de portée de cette session.

\textbf{Point crucial, à ne pas manquer} : ce diagnostic négatif est
\emph{spécifique à la version globale}. Avec des matrices globales,
CHAQUE produit matrice-vecteur touche potentiellement tout le maillage
(les non-zéros de $K_{cc}$, $K_{cf}$, $K_{fc}$ couvrent le domaine
entier), donc il n'y a \textbf{aucune localité spatiale à exploiter} --
paralléliser une seule opération globale, c'est distribuer un travail
déjà minuscule sur des cœurs qui n'ont individuellement presque rien à
faire chacun, d'où le surcoût dominant. La restriction (\S
\ref{sec:restrictedattempts}) n'est donc pas une alternative à la
parallélisation : c'est son \textbf{préalable}. Si l'approche restreinte
(coût $O(\text{taille de bande})$ par niveau) était rendue exacte, les
différentes bandes -- ou différentes régions du maillage traitées à un
même niveau de récursion -- deviendraient des unités de travail
\emph{spatialement localisées et mutuellement indépendantes}, exactement
le grain de parallélisme qui manque ici. C'est \emph{cette}
combinaison (restriction d'abord, parallélisation ensuite -- pas
l'inverse) qui donnerait un vrai gain, mais elle suppose d'avoir
résolu le résidu $\sim$15--17$\times$ documenté en
\S\ref{sec:restrictedattempts}--\ref{sec:diagn2}, ce qui n'a pas été
possible ce soir.

\section{État final de la session}

\begin{itemize}
  \item \textbf{Acquis} : algorithme multi-niveaux exact (matrices
  globales), validé à la fois par précision machine sur cas contrôlés
  et par correspondance à 9 chiffres avec la solution analytique sur un
  cas réel ($N=2$) ; recalibrage CFL donnant un gain $\sim$25$\times$
  confirmé et indépendant.
  \item \textbf{Non résolu} : aucune optimisation algorithmique
  (restriction, parallélisation) n'a amélioré le compromis
  précision/vitesse de la version globale exacte, malgré six tentatives
  documentées au total dans ce rapport.
  \item \textbf{En cours} : calcul de $N=3$ avec la version globale
  exacte + CFL recalibré, lancé en arrière-plan (ETA estimée plusieurs
  heures).
  \item \textbf{Résultat scientifique acquis} : 3 points de convergence
  ($N=0,1,2$) pour $k=3$, ordre $\approx4{,}60$, cohérent avec le
  résultat obtenu en début de session par une méthode indépendante.
\end{itemize}

\section{Le multiniveau réduit-il vraiment le coût ? Une vérification directe}
\label{sec:l2vsl5}

Question testée directement : sur le \emph{même} maillage $N=2$, avec le
\emph{même} algorithme exact (matrices globales), un découpage forcé à
$L=2$ (2-niveaux classique, sans niveaux intermédiaires) coûte-t-il plus
cher qu'un découpage automatique à $L=5$ ?

\begin{center}
\begin{tabular}{@{}lcc@{}}
\toprule
Configuration & Temps ($N=2$) & Erreur $L_2$ \\
\midrule
$L=2$ (forcé) & 104,98\,s & $1{,}144980932\times10^{-5}$ \\
$L=5$ (automatique) & 96--111\,s & $1{,}144978339\times10^{-5}$ \\
\bottomrule
\end{tabular}
\end{center}

Les deux configurations donnent un temps \textbf{quasi identique} (dans
le bruit de mesure) et la même erreur (correcte). \textbf{Explication} :
dans l'implémentation actuelle, le calcul \og propre \fg{} de chaque
niveau (\texttt{erk\_weight\_LTS\_coarse\_global}) utilise les matrices
\emph{globales} -- coût $O(n_{\text{dof}})$ à \emph{chaque} appel, quel
que soit le niveau. Le nombre total d'appels à ce calcul, sommé sur tous
les niveaux ancêtres $i=0,\dots,L-2$, vaut
\[
  \sum_{i=0}^{L-2} \prod_{j<i} p_j \;=\; p_0 + p_0 p_1 + p_0 p_1 p_2 + \cdots + p_0 p_1 \cdots p_{L-3} ,
\]
une somme de type géométrique \textbf{dominée par son dernier terme} --
qui est du même ordre de grandeur que $p_{\text{global}}$, la fréquence
d'appel unique du niveau coarse dans un schéma 2-niveaux classique
(pour $N=2$ : $p_0{=}4, p_0p_1{=}16, p_0p_1p_2{=}64, p_0p_1p_2p_3{=}256$,
somme $=340$, à comparer à $p_{\text{global}}{=}1024$ pour $L=2$ --
un facteur $\sim$3, pas plusieurs ordres de grandeur). Ajouter des
niveaux intermédiaires \emph{répartit} le même volume total de travail
sur plus d'appels à des fréquences différentes, sans le réduire
fondamentalement -- \textbf{le multiniveau ne devient réellement
avantageux que combiné à la restriction} (coût $O(\text{taille de
bande})$ au lieu de $O(n_{\text{dof}})$ par niveau), ce qui est
précisément ce que \S\ref{sec:restrictedattempts}--\ref{sec:moreattempts}
ont essayé d'obtenir sans y parvenir de façon exacte.

\section{Diagnostic fin : où et comment grandit l'écart restreint/exact}
\label{sec:diagn2}

Pour comprendre \emph{précisément} où l'écart entre version restreinte
(avec le correctif \texttt{advance\_uncovered\_full}, halo à 8 anneaux)
et version globale exacte se forme, un outil de diagnostic dédié
(\texttt{ERK4\_MLTS\_Lshape\_DiagN2\_test.hpp}) a été construit : il
fait tourner \emph{les deux} algorithmes en parallèle, à partir du
\emph{même} état initial, sur le \emph{vrai} maillage $N=2$ ($L=5$,
$p_{\text{global}}=1024$), pour un nombre variable de pas de temps
macro \texttt{nt}, puis compare bande par bande le degré de liberté où
l'écart est maximal.

\begin{center}
\begin{tabular}{@{}lccccc@{}}
\toprule
\texttt{nt} & band0 & band1 & band2 & band3 & band4 (terminal) \\
\midrule
1  & $1{,}17\times10^{-6}$ & $2{,}23\times10^{-6}$ & $2{,}13\times10^{-6}$ & $3{,}44\times10^{-6}$ & $7{,}71\times10^{-10}$ \\
2  & $2{,}14\times10^{-6}$ & $3{,}69\times10^{-6}$ & $2{,}49\times10^{-6}$ & $3{,}99\times10^{-6}$ & $5{,}45\times10^{-9}$ \\
3  & $2{,}96\times10^{-6}$ & $4{,}56\times10^{-6}$ & $2{,}45\times10^{-6}$ & $3{,}32\times10^{-6}$ & $2{,}00\times10^{-8}$ \\
5  & $4{,}25\times10^{-6}$ & $5{,}07\times10^{-6}$ & $3{,}19\times10^{-6}$ & $2{,}27\times10^{-6}$ & $4{,}85\times10^{-8}$ \\
30 & $1{,}04\times10^{-5}$ & $1{,}59\times10^{-5}$ & $1{,}21\times10^{-5}$ & $1{,}85\times10^{-5}$ & $1{,}90\times10^{-5}$ \\
\bottomrule
\end{tabular}
\end{center}
(valeurs : max$|x_{\text{restreint}} - x_{\text{global}}|$ sur les
degrés de liberté \emph{cellule} propres de chaque bande ; ce test
utilise \texttt{cfl\_factor}$=0{,}002$ par défaut, donc \texttt{nt}
petits ne représentent qu'une fraction du temps final simulé --
\texttt{cfl\_factor}$=0{,}002$ correspond à \texttt{nt}$=283$ pour
atteindre $t_f=0{,}05$.)

\textbf{Observations} :
\begin{itemize}
  \item L'écart après \emph{un seul} pas de temps macro est minuscule
  ($\sim10^{-6}$ à $10^{-10}$) -- la restriction n'introduit donc
  \emph{pas} d'erreur significative dans un pas isolé. Le bug (s'il y
  en a un) doit être un effet qui s'accumule sur \emph{plusieurs} pas.
  \item La croissance de \texttt{nt}$=1$ à \texttt{nt}$=30$ est
  \textbf{sous-linéaire} (facteur $\sim$2,45 pour un facteur 6 en
  \texttt{nt}, entre \texttt{nt}$=5$ et \texttt{nt}$=30$), \emph{pas}
  accélérée -- ce qui \textbf{contredit} l'hypothèse d'une instabilité
  qui s'aggrave avec le temps.
  \item Extrapoler cette tendance sous-linéaire jusqu'à
  \texttt{nt}$=283$ (le run complet à \texttt{cfl\_factor}$=0{,}002$)
  donne une estimation grossière de $\sim3$--$5\times10^{-5}$ --
  \textbf{loin} de l'écart final réellement observé entre la version
  restreinte et la solution analytique dans le run de production complet
  ($1{,}307\times10^{-4}$ à \texttt{cfl\_factor}$=0{,}002$,
  $1{,}535\times10^{-4}$ à \texttt{cfl\_factor}$=0{,}05$).
  \item \textbf{Ce résultat est non concluant, pas négatif} : soit il y
  a un effet de seuil/discontinuité entre \texttt{nt}$=30$ et
  \texttt{nt}$=283$ nécessitant plus de temps de calcul pour être
  caractérisé, soit la comparaison \og restreint vs.\ trajectoire
  globale identique \fg{} ne capture pas exactement le même phénomène
  que la comparaison finale \og restreint vs.\ solution analytique
  $\varphi(r,\theta)$ \fg{} utilisée dans les runs de production. Cette
  piste reste ouverte.
\end{itemize}

\section{Toutes les tentatives, en un coup d'œil}

\begin{center}
\begin{tabular}{@{}p{5.2cm}p{4cm}p{4.8cm}@{}}
\toprule
Tentative & Résultat & Section \\
\midrule
Correctif halo (\texttt{advance\_uncovered\_full}) & Marche sur petit cas contrôlé (ratio $\to$1,0), résidu $\sim$15--17$\times$ sur $N{=}2$ réel & \ref{sec:restrictedattempts} \\
Terme croisé \texttt{ancestor\_w} (Almquist-Mehlin) & Pire ($1{,}92\times10^{-4}$ vs $1{,}54\times10^{-4}$) & \ref{sec:moreattempts} \\
Micro-opt.\ (calcul gaspillé + scan redondant) & Correct, mais aucun gain de temps mesurable & \ref{sec:moreattempts} \\
OpenMP (parallélisme par appel) & Beaucoup plus lent (surcoût de synchronisation) & \ref{sec:parallel} \\
Hybride (ancêtres exacts, terminal seul restreint) & Pire sur les deux plans (précision \emph{et} vitesse) & \ref{sec:moreattempts} \\
Comparaison $L{=}2$ vs $L{=}5$ (matrices globales) & Coût quasi identique -- confirme que la restriction, pas le multiniveau seul, est la clé & \ref{sec:l2vsl5} \\
Diagnostic croissance par bande & Sous-linéaire jusqu'à \texttt{nt}$=30$, n'explique pas l'écart final -- non concluant & \ref{sec:diagn2} \\
\bottomrule
\end{tabular}
\end{center}

\section{Mise à jour finale : le bug d'ordre trouvé, un résidu persiste}
\label{sec:orderingfix}

Après l'écriture des sections précédentes, une nouvelle piste a été
identifiée et validée : un \textbf{bug d'ordre des opérations}, réel et
partiellement corrigé.

\textbf{Le bug} : dans le driver restreint, \texttt{advance\_uncovered\_full}
(la mise à jour directe en forme fermée des degrés de liberté propres
d'une bande $i$) était appelée \emph{avant} de récurser dans la bande
$i+1$. Or dans l'algorithme global exact, \texttt{xn} n'est \emph{jamais}
écrit par un niveau ancêtre -- seul le niveau terminal écrit, et tout
calcul de \texttt{own\_w} est une lecture pure de l'état tel qu'il était
\emph{avant} cet intervalle macro. Écrire dans \texttt{xn} avant de
récurser permettait au calcul de \texttt{own\_w\_\{i+1\}} de la bande
$i+1$ de lire un état déjà contaminé par l'écriture directe de la bande
$i$ -- un bug d'ordre absent de la version globale par construction.

\textbf{Le correctif} : déplacer l'appel à
\texttt{erk\_weight\_LTS\_coarse\_advance\_uncovered\_full} pour qu'il
ait lieu \emph{après} l'appel récursif \texttt{recurse(i+1, ...)}, pas
avant, à chaque niveau de la récursion.

\textbf{Résultat sur $N=2$ réel} (\texttt{cfl\_factor}$=0{,}05$) :
\begin{center}
\begin{tabular}{@{}lcc@{}}
\toprule
& Erreur $L_2$ & Temps \\
\midrule
Avant correctif d'ordre & $1{,}535\times10^{-4}$ & 88,8\,s \\
\textbf{Après correctif d'ordre} & $\mathbf{8{,}35\times10^{-5}}$ & \textbf{40,9\,s} \\
Référence (exacte) & $1{,}145\times10^{-5}$ & -- \\
\bottomrule
\end{tabular}
\end{center}
Une réduction d'erreur de \textbf{45\%} et un gain de vitesse de
\textbf{2$\times$}, sans aucune régression sur les cas contrôlés
($N=0,1$ restent exacts).

\textbf{Hypothèses testées ensuite pour fermer le résidu restant
($\sim$7,3$\times$ la référence), toutes négatives} :
\begin{itemize}
  \item \textbf{Largeur du halo} (8 $\to$ 20 anneaux) : quasi aucun effet
  ($7{,}96\times10^{-5}$ vs $8{,}35\times10^{-5}$) -- confirme que le
  halo n'est pas le facteur limitant (cohérent avec la précision machine
  déjà mesurée pour \texttt{own\_w\_i}, \S\ref{sec:diagn2}).
  \item \textbf{Supprimer l'avance directe des bandes intermédiaires}
  (ne garder que band0) : \textbf{explosion catastrophique}
  ($L_2\sim10^{24}$) -- confirme que chaque bande non-terminale a
  \emph{absolument besoin} de sa propre avance directe ; ce n'est pas
  redondant avec la fuite du terminal.
\end{itemize}

\textbf{État du code à la fin de la session} : le fichier
\texttt{ERK4\_MLTS\_Lshape\_RestrictedExact\_conv\_test.hpp} contient le
correctif d'ordre (actif), \texttt{halo\_rings}$=8$ (restauré, la
largeur 20 n'apportait rien), et l'avance directe active pour
\emph{toutes} les bandes non-terminales. C'est le \textbf{meilleur état
validé} à ce jour : rapide, sans régression sur les cas contrôlés, mais
avec un résidu $\sim$7$\times$ la référence sur $N=2$ réel, cause
encore non identifiée.

\textbf{Piste ouverte pour la prochaine session} : le résidu ne vient ni
de la précision de \texttt{own\_w\_i} (prouvée exacte), ni du halo, ni
d'une bande superflue -- il reste soit un second bug d'ordre/lecture
plus subtil (peut-être entre bandes \emph{sœurs}, pas seulement
ancêtre-descendant), soit une interaction fine entre l'accumulation
\texttt{combined = shifted + own\_w\_i} et le moment précis où chaque
bande \og voit \fg{} les contributions des autres. Le diagnostic
\texttt{ERK4\_MLTS\_Lshape\_DiagN2\_test.hpp} (comparaison directe,
bande par bande, contre la version globale exacte, sur le vrai maillage
$N=2$) reste l'outil le plus efficace pour investiguer -- il isole le
problème en quelques secondes à quelques minutes de calcul, sans
attendre un run complet.

\section{Nouvelle session : bug de régression identifié par comparaison
inter-fichiers, correctif stable trouvé}
\label{sec:regressionfix}

Cette section documente une nouvelle session d'analyse, menée en
s'appuyant explicitement sur une relecture ciblée de Almquist--Mehlin
(CRC Preprint 2016/16) et Diaz--Grote (CMAME 2015), pour auditer
\texttt{ERK4\_MLTS\_Lshape\_RestrictedExact\_conv\_test.hpp} --- le
\og meilleur état validé \fg{} hérité de la session précédente
(\S\ref{sec:orderingfix}), avec son résidu $\sim$7,3$\times$ non
expliqué sur $N=2$ réel.

\subsection{Diagnostic : une régression trouvée par grep, pas par
raisonnement seul}
\label{sec:regressiondiag}

Un grep de tous les appels à \texttt{erk\_weight\_LTS\_coarse\_advance\_-
uncovered[\_full]} dans le dossier \texttt{prototypes/LTS/} révèle un
pattern net :

\begin{center}
\begin{tabular}{@{}p{7.5cm}p{3cm}p{2.2cm}@{}}
\toprule
Fichier & Argument passé & Statut \\
\midrule
\texttt{ERK4\_MLTS\_L2\_validation\_test.hpp:238} & \texttt{combined} & correct \\
\texttt{ERK4\_MLTS\_L3\_validation\_test.hpp:295} & \texttt{combined} & correct \\
\texttt{ERK4\_MLTS\_L5small\_validation\_test.hpp:295} & \texttt{combined} & correct \\
\textbf{\texttt{...RestrictedExact\_conv\_test.hpp:354}} & \textbf{\texttt{own\_w\_i}} & \textbf{régression} \\
\texttt{ERK4\_MLTS\_RestrictedGD\_test.hpp:399} & \texttt{own\_w\_i} & même régression \\
\texttt{ERK4\_MLTS\_Lshape\_DiagN2\_test.hpp:387} & \texttt{own\_w\_i} & même régression \\
\bottomrule
\end{tabular}
\end{center}

Les fichiers \texttt{*\_validation\_test.hpp} (ceux qui avaient validé
la récursion générique à précision quasi machine, \S\ref{sec:mlts}) font
tous \texttt{combined = ancestor\_w + own\_w} avant d'appeler
\texttt{advance\_uncovered}. Les fichiers plus récents avaient régressé
vers \texttt{own\_w\_i} seul --- vraisemblablement en généralisant depuis
les variantes \texttt{*\_v2design} où seule \texttt{blocks[0]} (la
racine, sans ancêtre) était traitée, cas où \texttt{own\_w == combined}
par coïncidence, ce qui masquait la régression.

\textbf{Justification théorique (Almquist--Mehlin, éq. 18--22)} :
\texttt{erk\_weight\_LTS\_coarse\_v2} est appelé avec
\texttt{ancestor\_w = nullptr} par défaut ; son propre commentaire
indique que la vraie ODE locale d'une bande est
$\dot x = B(x) + M_c^{-1}F(t) + \text{ancestor\_w}(t)$, donc
\texttt{own\_w\_i} seul ne représente que les deux premiers termes.
L'équation (21) de l'article confirme la structure additive attendue :
$F_L(t^{[L]}) = \dots + \sum_{\ell=0}^{L-1}(P_\ell-P_{\ell+1})\,
r_\ell(\dots)$ --- chaque bande ancêtre $\ell$ doit contribuer à ses
\emph{propres} degrés de liberté (sélectionnés par $P_\ell-P_{\ell+1}$),
en plus de sa propre polynôme. Pour la bande 0 (racine), \texttt{shifted}
est toujours nul (pas d'ancêtre) : \texttt{combined == own\_w\_i}
exactement, ce qui explique pourquoi la magnitude de band0 était
toujours correcte alors que les bandes intermédiaires ($i\geq1$, où
\texttt{shifted}$\neq0$) étaient silencieusement sous-forcées ---
cohérent avec le résidu croissant avec $L$ et non avec $p_{global}$
(\S\ref{sec:lgrowth}).

\subsection{Premier correctif (addition naïve) : rejeté, cause une
explosion}
\label{sec:naivefail}

Un premier correctif, minimal, remplace simplement \texttt{own\_w\_i}
par \texttt{combined = shifted + own\_w\_i} (déjà calculé, déjà en
scope) dans l'appel à \texttt{advance\_uncovered\_full}. Testé sur
$N=1$ (trivial, $L=2$, sans ancêtre donc sans effet) puis sur $N=2$
réel ($L=5$) : \textbf{explosion catastrophique},
$L_2\sim8{,}9\times10^{19}$.

\textbf{Diagnostic} : contrairement à ce que suggère la référence
globale exacte (\S\ref{sec:globalfix}, qui utilise elle aussi une
addition naïve \texttt{combined = shifted + own\_w\_i} sans jamais
planter), la fonction \texttt{advance\_uncovered\_full} n'utilise PAS
\texttt{combined} comme un simple terme de forçage échantillonné
ponctuellement (ce que fait RK4 aux points de quadrature) : elle
l'utilise comme les coefficients de Taylor \emph{exacts} d'une
intégrale de Simpson fermée,
$\Delta x = w_0\,dt + w_1\,dt^2/2 + w_2\,dt^3/6 + w_3\,dt^4/24$. Si
\texttt{combined} n'est pas le \emph{vrai} polynôme cubique de la
dynamique couplée (parce que l'addition naïve ignore que le forçage
ancêtre $g(t)$ doit lui-même être re-propagé à travers l'opérateur $B$
--- termes croisés $Bg_0$, $B^2g_0$, $Bg_1$, exactement comme le fait
$F(t)$ dans le même calcul), l'intégrale de Simpson amplifie cette
erreur d'ordre supérieur au lieu de simplement l'échantillonner ---
un mécanisme d'instabilité spécifique au raccourci en forme fermée,
absent de la version globale (qui ne fait jamais d'intégrale directe :
seul le niveau terminal écrit \texttt{xn}, via RK4, jamais un
quadrature en forme close).

\subsection{Correctif retenu : termes croisés via \texttt{ancestor\_w},
pas d'addition séparée}
\label{sec:crossfermfix}

Le correctif stable fait transiter \texttt{shifted} \emph{dans}
\texttt{erk\_weight\_LTS\_coarse\_v2} via son paramètre
\texttt{ancestor\_w} (déjà implémenté et documenté dans
\texttt{erk\_coupling\_hho\_scheme.hpp}, mais jamais branché
correctement jusqu'ici) :

\begin{lstlisting}[language=C++]
erk_an.erk_weight_LTS_coarse_v2(xn, blocks[i], own_w_i,
    Fn, Fn12, Fn1, dtau_i, &shifted);   // ancestor_w passe en 8e argument
recurse(i + 1, t_start + tm, dtau_i, xn, own_w_i);   // deja complet
erk_an.erk_weight_LTS_coarse_advance_uncovered_full(
    xn, blocks[i], own_w_i, blocks[L-1].active_c, blocks[L-1].active_f, dtau_i);
\end{lstlisting}

\texttt{own\_w\_i} devient alors directement le polynôme total
correct (avec termes croisés), utilisé \emph{tel quel} --- à la fois
comme forçage transmis à la récursion et pour l'avance en forme
fermée --- sans aucune addition séparée. C'est important : une
tentative antérieure de la session précédente (\S\ref{sec:moreattempts},
point 1, \og terme croisé ancestor\_w \fg{}) avait déjà essayé de
brancher \texttt{ancestor\_w} et trouvé un résultat \emph{pire}
($1{,}92\times10^{-4}$ contre $1{,}54\times10^{-4}$) --- diagnostic
confirmé après coup : cette tentative gardait \emph{en plus} une
addition \texttt{combined = shifted + own\_w\_i} après l'appel, ce
qui double-comptait le terme d'ordre 0 de \texttt{shifted} (déjà inclus
dans \texttt{own\_w\_i} via \texttt{MinvF0\_tot}). Le correctif retenu
ici supprime cette addition redondante.

\subsection{Résultats}

\begin{center}
\begin{tabular}{@{}lccc@{}}
\toprule
Configuration ($N=2$, $L=5$, $p_{global}=1024$, \texttt{cfl}$=0{,}05$) & Erreur $L_2$ & Ratio/réf. & Temps \\
\midrule
Avant cette session (bug \texttt{own\_w\_i}, ordre déjà corrigé) & $8{,}35\times10^{-5}$ & $\times7{,}3$ & 40,9 s \\
Addition naïve \texttt{combined} (rejetée) & $8{,}9\times10^{19}$ & explosion & -- \\
\textbf{\texttt{ancestor\_w}, sans addition séparée (retenu)} & $\mathbf{6{,}74\times10^{-5}}$ & $\mathbf{\times5{,}9}$ & \textbf{43,9 s} \\
Référence (globale exacte) & $1{,}145\times10^{-5}$ & $\times1$ & $\sim$100 s \\
\bottomrule
\end{tabular}
\end{center}

Amélioration nette et stable (résidu $\times7{,}3\to\times5{,}9$, aucune
instabilité), sans régression sur $N=0,1$ (exacts par construction, un
seul niveau ou zéro ancêtre). Correctif propagé également à
\texttt{ERK4\_MLTS\_RestrictedGD\_test.hpp} et
\texttt{ERK4\_MLTS\_Lshape\_DiagN2\_test.hpp}, qui partageaient soit la
même régression, soit le même risque d'explosion.

\textbf{Élargissement du halo re-testé (8$\to$20 anneaux)} : confirmé
nuisible ($1{,}52\times10^{-3}$, bien pire, et 84 s au lieu de 44 s) ---
cohérent avec le rapport précédent, cette piste reste définitivement
écartée. Le résidu $\times5{,}9$ restant n'est donc pas expliqué par la
troncature du halo ; sa cause exacte demeure ouverte pour une future
session (candidats : la nature intrinsèquement différente d'une
intégrale de Simpson en forme fermée par bande, appliquée à la
fréquence \emph{propre} de chaque bande, comparée à l'échantillonnage
RK4 du niveau terminal à la fréquence \emph{la plus fine} dans la
version globale).

\subsection{Gain de performance mis en dur}

Le \texttt{cfl\_factor} par défaut des deux drivers
(\texttt{...RestrictedExact\_conv\_test.hpp} et
\texttt{...GlobalExact\_conv\_test.hpp}) est passé de $0{,}002$ à
$0{,}05$ --- gain $\times25$ déjà validé (\S\ref{sec:cflfix}), plus
besoin de la variable d'environnement \texttt{MLTS\_CFL\_FACTOR} pour
en bénéficier.

Avec ce correctif et le CFL recalibré, le point $N=2$ complet
(restreint, corrigé) s'exécute en $\sim$44 s, contre $\sim$100 s pour
la version globale exacte à $\texttt{cfl}=0{,}05$ --- un gain
$\times2{,}2$ tout en étant \emph{plus précis} qu'avant ce correctif.
Un balayage $N=3$ ($L=8$ niveaux, $p_{global}=32768$, précédemment
estimé à plusieurs jours avec la version globale non recalibrée) a été
lancé pour évaluer si l'algorithme restreint corrigé rend enfin ce
point de maillage praticable.

\end{document}
